\chapter{Hệ điều hành mã nguồn mở và góc nhìn ATTT}
\label{ch:os}
\section{Phạm vi so sánh và kiến trúc kernel}
Tại thời điểm đối chiếu ngày 27/10/2026, ảnh cài đặt LTS mới nhất trên nguồn Canonical là Ubuntu 26.04.1 LTS. Đây là bản cập nhật điểm của nhánh 26.04, không phải phiên bản đang cài trên máy thực hành. Máy của nhóm trả về Ubuntu 24.04.4 LTS trong \code{/etc/os-release}. Windows 11 được dùng làm đối tượng so sánh kiến trúc và chính sách, không phải môi trường đã chạy website trong đợt kiểm thử này. Không suy diễn rằng mọi hành vi trên Ubuntu 26.04 đã được đo thực nghiệm từ máy 24.04.4. Nhánh 26.04 có hỗ trợ tiêu chuẩn đến năm 2031; các giai đoạn mở rộng có điều kiện dịch vụ/phạm vi gói riêng.\cite{ubuntuRelease,ubuntuCycle}

Linux tổ chức bộ nhớ đặc quyền kernel tách khỏi tiến trình ở user space. Ứng dụng như PHP hoặc trình duyệt sử dụng thư viện và system call để yêu cầu kernel mở file, tạo tiến trình hay gửi dữ liệu mạng. Điểm chuyển quyền này cần kiểm tra tham số, con trỏ và quyền của chủ thể yêu cầu; một đường dẫn HTTP không trực tiếp trao quyền kernel cho mã PHP. Linux thường được mô tả là kernel nguyên khối có module: nhiều dịch vụ hệ thống và driver cùng thực thi trong kernel space, trong khi module cho phép nạp thành phần theo nhu cầu mà không biên dịch lại toàn bộ kernel. Khả năng nạp module hữu ích cho phần cứng nhưng cũng tạo bề mặt đặc quyền cao, nên không cấp quyền quản trị máy cho tài khoản vận hành nội dung WordPress.\cite{linuxSyscall,linuxModules}

Windows 11 sử dụng dòng kiến trúc Windows NT, kết hợp kernel, executive và các driver. Ứng dụng chạy ở user mode, tương tác qua các API và dịch vụ hệ thống. Kernel và các thành phần executive phụ trách những chức năng như bộ nhớ, tiến trình, I/O và đối tượng; driver có thể chạy ở kernel mode hoặc user mode tùy loại. Lỗi trong thành phần kernel có thể ảnh hưởng rộng hơn một tiến trình ứng dụng. Cả hai hệ điều hành đều có ranh giới đặc quyền và cơ chế kiểm soát truy cập, vì vậy không nên kết luận Linux an toàn chỉ vì có module hoặc Windows kém an toàn chỉ vì tổ chức executive khác Linux.\cite{windowsModes,windowsExecutive}

Trong kiến trúc website của nhóm, Nginx tiếp nhận request rồi chuyển PHP sang PHP-FPM; MySQL là dịch vụ dữ liệu riêng. Các tiến trình đó không cần chạy bằng root để xử lý bài ảnh. Nếu một plugin bị khai thác ở mức PHP, phạm vi thiệt hại chịu ảnh hưởng bởi quyền filesystem, thông tin DB mà PHP đọc được và các dịch vụ có thể truy cập. Ranh giới kernel vẫn quan trọng, nhưng không thay thế việc kiểm tra nonce, capability và MIME trong ứng dụng. Tách tài khoản vận hành, cập nhật driver/kernel và hạn chế module ngoài nguồn tin cậy là các lớp kiểm soát bổ trợ.

\section{Phân quyền Linux và Windows}
Trên Linux, UID biểu diễn người dùng; GID cùng các nhóm bổ sung tạo bối cảnh nhóm. Mỗi inode có owner/group và các bit quyền cho owner, group, others. Với file thường, \code{r} cho đọc dữ liệu, \code{w} cho sửa nội dung, \code{x} cho thực thi. Với thư mục, \code{r} cho liệt kê tên, \code{x} cho đi xuyên/tra cứu các mục con, còn \code{w} kết hợp với quyền đi xuyên cho phép tạo hoặc xóa mục trong thư mục. Vì vậy, một file chỉ đọc vẫn có thể bị xóa bởi chủ thể có quyền thích hợp trên thư mục cha; chỉ đổi file về 644 không đủ bảo vệ nếu thư mục chứa nó cho người không tin cậy ghi.

ACL mở rộng cho phép cấp quyền cụ thể cho nhiều UID/GID thay vì chỉ ba lớp truyền thống. Mask ACL giới hạn quyền hiệu lực của các mục liên quan; cần kiểm tra \code{getfacl} khi kết quả truy cập không khớp phép đọc đơn giản từ \code{ls -l}. Default ACL trên thư mục có thể ảnh hưởng quyền của file mới, còn umask ảnh hưởng quyền lúc tạo. Chmod một lần không bảo đảm mọi ảnh được plugin sinh về sau đều có cùng quyền. Trong triển khai LocalWP phải đối chiếu cả owner của PHP-FPM, Nginx và đường đi qua các thư mục cha.\cite{linuxAcl}

Windows dùng SID định danh người dùng/nhóm. Access token gắn với tiến trình hoặc luồng mang SID, nhóm và privilege. Security descriptor của một đối tượng có thể chứa DACL và SACL. DACL gồm ACE cho phép hoặc từ chối truy cập; SACL xác định những trường hợp cần audit, không phải danh sách cấp quyền bổ sung. Inheritance truyền các ACE được đánh dấu xuống đối tượng con. DACL rỗng và không có DACL khác nhau: rỗng không cấp quyền, còn trường hợp không có DACL có thể cho phép truy cập rộng. Cần dùng công cụ quản trị/API phù hợp thay vì diễn giải tất cả bằng mô hình số chmod.\cite{windowsTokens,windowsAcl}

\begin{table}[H]\centering\small
\caption{Ý nghĩa quyền áp dụng trên website của nhóm}
\begin{tabularx}{\linewidth}{lYY}\toprule Đối tượng & Quyền và điều kiện & Ý nghĩa ATTT\\\midrule
Thư mục ứng dụng & 755, owner phù hợp tiến trình & Owner quản lý; nhóm/others đi xuyên và đọc, không tự ghi.\\
File PHP/CSS/JS & 644 & PHP cần đọc mã; không cần bit thực thi cho từng file PHP qua FPM.\\
\code{wp-config.php} & 400 thực tế; 440 nếu dùng group phù hợp & Hạn chế đọc secret; phải kiểm tra PHP vẫn đọc được.\\
Dữ liệu Wordfence & Có file 600 & Giữ mức chặt hơn 644; không nới quyền máy móc.\\\bottomrule\end{tabularx}\label{tab:permissions}
\end{table}
Nguyên tắc chung là quyền tối thiểu: người gửi ảnh chỉ có quyền nội dung cần thiết, còn tiến trình máy chủ chỉ ghi vào vùng vận hành cần thiết. Chmod 777 làm mất ranh giới ghi; chmod đệ quy toàn home/Git có thể phá ứng dụng khác. Với LocalWP, các dịch vụ chạy cùng tài khoản desktop nên mô hình cách ly còn hạn chế; kết quả quyền local không tự là cấu hình production tối ưu.

\section{Quản lý bản vá}
APT quản lý gói, dependency và phiên bản qua kho cấu hình trên Ubuntu. Thông tin kho có chữ ký giúp xác nhận nguồn và tính toàn vẹn của metadata; cần giữ key/kho đáng tin, kiểm tra nguồn gói và tránh thêm PPA không cần thiết. \code{unattended-upgrades} có thể tự áp dụng các cập nhật được chính sách cho phép. Quản trị viên vẫn cần đọc log, theo dõi cập nhật bị giữ lại và lên lịch reboot khi thành phần đang chạy chưa dùng bản mới. Bản vá cài trên đĩa chưa đồng nghĩa kernel hoặc dịch vụ đang chạy đã được thay thế.\cite{ubuntuUpdates}

Windows Update cung cấp cơ chế phân phối cập nhật hệ điều hành. Các chính sách cho máy tổ chức có thể quản lý thời điểm nhận bản cập nhật, vòng triển khai và trải nghiệm khởi động lại qua Group Policy hoặc MDM, tùy edition. Phương pháp triển khai theo nhóm thử nghiệm rồi mở rộng giúp phát hiện lỗi tương thích trước khi ảnh hưởng toàn hệ thống. Hoãn cập nhật cần thời hạn và giám sát; hoãn vô thời hạn biến một biện pháp kiểm soát ổn định thành tồn đọng lỗ hổng.\cite{windowsUpdates}

Vòng đời Windows 11 phụ thuộc edition và bản feature: nguồn Microsoft công bố chu kỳ hỗ trợ thông thường 24 tháng cho Home/Pro và 36 tháng cho Enterprise/Education; LTSC có lịch riêng. Vì vậy phải tra đúng phiên bản thực tế thay vì nói mọi Windows 11 cùng ngày hết hỗ trợ. Sau cài bản vá, cần xác nhận build đang chạy, reboot khi cần và chuẩn bị phương án gỡ bản cập nhật/phục hồi đã thử, tương tự việc kiểm tra kernel/dịch vụ sau cập nhật trên Ubuntu.\cite{windowsLifecycle}

Quy trình phù hợp với website của nhóm gồm kiểm kê phiên bản; theo dõi thông báo chính thức; sao lưu file, DB và template LocalWP; thử trên clone; chạy luồng upload--duyệt--hiển thị; sau đó mới cập nhật môi trường sử dụng. Kế hoạch rollback phải xác định rõ phục hồi riêng gói, template hay DB. Khi schema thay đổi, khôi phục chỉ code có thể không tương thích DB đã nâng cấp. Khi hỏng rewrite WebP, phục hồi đúng template hiệu quả hơn phục hồi toàn bộ database.

Vá Ubuntu không tự vá runtime nằm trong LocalWP: PHP/Nginx/MySQL được LocalWP quản lý riêng. WordPress core, Astra và các plugin cũng có vòng phát hành độc lập. Mỗi lớp cần kiểm kê và kiểm tra lại. Trong đợt này không nâng hàng loạt phiên bản để làm đẹp kết quả scanner. Công cụ nhận diện phiên bản chỉ cung cấp đầu vào đánh giá; bản vá cần đối chiếu advisory và điều kiện ảnh hưởng, không suy CVE từ một chuỗi số.

\section{GPL, giấy phép thương mại và nghĩa vụ phân phối}
Linux kernel được công bố theo GPL-2.0-only ở cấp tổng thể, có ngoại lệ syscall và quy tắc cho từng file tương thích. Đây không phải tuyên bố mọi gói trong Ubuntu đều cùng GPL: thư viện, công cụ và tài nguyên có thể mang GPL các phiên bản khác, LGPL, BSD, MIT hoặc giấy phép khác. Khi phân phối, cần xem đúng giấy phép của thành phần và hình thức phân phối, giữ thông báo tác quyền, điều kiện kèm theo và đáp ứng nghĩa vụ cung cấp mã nguồn tương ứng khi áp dụng. Việc chương trình user space gọi syscall không tự làm mọi chương trình đó trở thành GPL.\cite{linuxKernel,ubuntuLicensing}

Windows 11 được cung cấp theo điều khoản phần mềm sở hữu độc quyền; quyền cài đặt, chuyển giao và sử dụng phụ thuộc edition, kênh OEM/retail/volume và hợp đồng tương ứng. Người dùng không mặc nhiên có quyền sửa rồi tái phân phối mã nguồn Windows như đối với Linux kernel. Phần mềm mã mở có thể được bán hoặc cung cấp dịch vụ thương mại; thu phí không làm nó tự mất tính mã mở. Ngược lại, có thể tải một phần mềm không mất tiền không đồng nghĩa người dùng được quyền sửa và phân phối mã nguồn.\cite{windowsTerms}

Đối với Lăng Kính, file LICENSE gốc MIT mang thông báo tác quyền hiện hữu được giữ nguyên. Child theme và module bảo mật khai báo GPL-2.0-or-later ở phạm vi thành phần của chúng; không áp MIT gốc để đổi giấy phép Astra/WordPress/plugin. Image Watermark bản đã cài khai báo MIT; Wordfence khai báo GPL-3.0-or-later. Giấy phép mã không tự cấp quyền dùng ảnh, font hay nhận diện thương hiệu. Bàn giao học phần cần kèm bảng kiểm kê và giới hạn quyền ảnh chưa được xác nhận, thay vì cam kết mọi tài sản đều do nhóm sở hữu.

\section{Auditability và chuỗi cung ứng}
Khả năng đọc source cho phép kiểm tra đường đi dữ liệu, rà dependency, tái hiện lỗi và đề xuất bản sửa. Lỗi Pass Thru của dự án minh họa ích lợi này: có thể đối chiếu phần đổi URL với danh sách đường dẫn mà loader thực sự cho phép. Tuy nhiên source công khai chỉ tạo điều kiện audit; chất lượng còn phụ thuộc năng lực người kiểm tra, mức độ kiểm thử và việc source có đúng với artifact đang chạy hay không.

Một chuỗi xác minh hữu ích đi từ nguồn tải chính thức, checksum/chữ ký và phiên bản gói, đến hash file thực thi, cấu hình đang có hiệu lực và hành vi HTTP. Hash source/live giúp phát hiện chép nhầm child theme; không tự chứng minh mã không có lỗ hổng. Tương tự, chữ ký hợp lệ chứng minh nguồn/kênh phát hành theo mô hình tin cậy, không chứng minh mọi logic trong bản phát hành đều an toàn. Kho dependency, plugin và dịch vụ cập nhật vẫn là những điểm cần giám sát.

Mã đóng hạn chế khả năng rà soát độc lập toàn bộ source, nhưng không tự đồng nghĩa có mã độc. Có thể đánh giá bằng chữ ký nhà phát hành, chính sách cập nhật, tài liệu kỹ thuật, log và kiểm thử hành vi. Kết luận ATTT cần kết hợp các bằng chứng, tránh dùng nhãn ``mã mở'' hoặc ``mã đóng'' như một chứng chỉ an toàn. Bản bàn giao vì vậy giữ source tự viết, cấu hình mẫu không bí mật, test có thể chạy lại và log đã giới hạn dữ liệu, đồng thời nêu rõ những phần chưa kiểm chứng.
